\documentclass[12pt]{article}
\usepackage{hyperref}

\usepackage{graphicx}

\usepackage{mathtools}
\DeclarePairedDelimiter\ceil{\lceil}{\rceil}
\DeclarePairedDelimiter\floor{\lfloor}{\rfloor}


\usepackage{amsmath}
\usepackage{amsfonts}
\usepackage{amssymb}

\newcommand{\Mod}[1]{\ (\mathrm{mod}\ #1)}


\usepackage{listings}
\lstset{language=C++,
                basicstyle=\ttfamily,
                keywordstyle=\color{blue}\ttfamily,
                stringstyle=\color{red}\ttfamily,
                commentstyle=\color{green}\ttfamily,
                morecomment=[l][\color{magenta}]{\#}
}

\usepackage{breqn}
\usepackage[]{algorithm2e}
\usepackage{physics}

\newcommand\fnurl[2]{%
  \href{#2}{#1}\footnote{\url{#2}}%
}

\newcommand{\Four}{\mathcal{F}}
\renewcommand{\(}{\left(}
\renewcommand{\)}{\right)}
%% \renewcommand{\{}{\left{}
%% \renewcommand{\}}{\right}}


\title{VASP notes for rocFFT}
\author{Malcolm Roberts, \texttt{malcolm.roberts@amd.com}}



\def\no#1{\nobreak{#1}} % stop equation breaking in dmath

\begin{document}
\maketitle

VASP, from a FFT library perspective, solves the nonlinear
Schrodinger equations,
\begin{dmath}
  \label{nls}
  i\partial_t\psi = -\frac{1}{2} \partial_x^2 \psi + \kappa \abs{\psi}^2\psi
\end{dmath}
using the pseudospectral method.  Equation~\eqref{nls} has a
third-order nonlinearity, but also has other forms which may involve
higher-order terms.  We note that $\psi$ is complex, and does not
exhibit Hermitian symmetry.

\section{The pseudospectral method and convolutions}

The pseudospectral method uses a Fourier transform to linearize the
spatial derivative, and transforms the nonlinear term into a
convolution.  From my understanding, the wave numbers for $\psi$ are
symmetric about the origin.  Let $\hat\psi$ denote the discrete
Fourier transform of $\psi$.  Then,
\begin{equation}
  \hat\psi = \left\{\hat\psi_{\mathbf{k}} \right\},
  \mathbf{k} = \(k_x, k_y, k_z\)
\end{equation}
where $k_x/k_0 \in \{-N_x +1, \dots, N_x -1\}$, and similarly for
$k_y$ and $k_z$ (and we note that $k_0$ may have different values for
different directions).  We refer to this type of data, where the
indices are symmetric about the origin, as centred.

The nonlinear term is mapped to a convolution:
\begin{equation}
  \abs{\psi}^2\psi
  \xrightarrow[\mathcal{F}]{}
  \widehat{\(\abs{\psi}^2\psi\)}_\ell
  =
  \sum_{a}\sum_{b}\sum_{c} \abs{\psi}_a \abs{\psi}_b \psi_c \delta_{\ell, a+b+c}.
\end{equation}
Direct computation of the convolution has complexity
$\mathcal{O}(N^3)$ and relatively poor numerical accuracy; this can be
improved by using an FFT-based method.  However, the resulting
convolution has extra terms from aliasing.  In order to recover the
original convolution, we need to remove these terms, ie dealias the
FFT-based convolution.

Phase-shift dealiasing and zero-padding are methods for dealiasing
convolutions; we consider zero-padding here. For binary convolutions
on 1D non-centred data, one must pad from length $N$ to length $2N-1$
or more.  For centred data, one needs to pad from length $2N$ to
length $3N$, which is referred to as $2/3$ padding.  Since the
convolution from equation~\eqref{nls} is ternary, we must pad from
length $2N$ to length $4N$ in order to fully dealias the convolution.

rocFFT is exploring adding convolution operations (and other spectral
operations, eg differentiation) to the rocFFT API.  In order to ensure
that this is useful, we are trying to understand use cases and
prototype API designs.

Questions:
\begin{itemize}
\item Is there ever a NLS with quadratic nonlinearity?
\item Is VASP (or other quantum simulators) always interested in
  third-order nonlinearities?
\item Do you use the rocFFT or hipFFT API for calling rocFFT?
\end{itemize}

\section{Problem sizes}

From our last discussion, VASP uses very small FFTs; $32^3$ was
mentioned.  This problem size actually fits in LDS/shared memory, so
it's possible to compute this with a single global-memory round-trip.
Currently, we compute this as a 2D kernel and a 1D kernel.  We may be
able to compute this as a single kernel, which would reduce
computation time and latency. We hope to eventually generate kernels
for all problem sizes that fit in LDS, but knowing useful problem sizes
would help prioritization.

Questions:
\begin{itemize}
\item Besides $32^3$, what problem sizes are important?
\end{itemize}

\section{Sparsity and multi-gpu transforms}

We are also scoping multi-gpu transforms, where we are hoping to add a
brick-based decomposition.  The single-process multi-gpu functional
implementation is scheduled for ROCm 6.0, and is available in the
public github repository.  A feature that works here as well is to
label certain bricks as sparse, ie as a region that is uniformly zero.
The proposed rocFFT convolution routines would handle zero-padding
internally, but other sparsity patterns can be handled by the
brick-based API.  It was mentioned that, for VASP, the non-zero data
is contained in the sphere inscribed in the cubic computational
domain, which seems like a useful application for sparsity.


Questions:
\begin{itemize}
\item What data decomposition can we expect from VASP?
\item Is the inscribed sphere the general sparsity pattern?
\end{itemize}

\section{Memory management}

rocFFT features an API for specifying work memory, which is specified
at the execution stage
(\lstinline{rocfft_execution_info_set_work_buffer}); this can help
reduce memory footprint by allowing buffer-reuse.  This API isn't
available in cuFFT, and not implemented in the hipFFT wrapper for the
rocFFT back-end.  Also, rocFFT makes single-direction plans with a
specific placement (ie in-place or out-of-place), so plans are more
targeted, which means that we can look at fewer cases, which also
helps reduce memory footprint.

We are also looking at adding functionality to allow the user to
prioritize minimizing memory footprint or minimizing execution time.

Questions:
\begin{itemize}
\item How important would work buffer specification be for hipFFT?
\item How high of a priority would optimizing for memory/time be for VASP?
\end{itemize}

\end{document}

% LocalWords:  VASP FFT pseudospectral nls nonlinearity linearize ie
% LocalWords:  dealias dealiasing rocFFT eg nonlinearities hipFFT LDS
% LocalWords:  FFTs gpu ROCm github rocfft cuFFT
